\documentclass[12pt]{article}
\usepackage{amsmath,amssymb,amsfonts}
\usepackage[margin=1in]{geometry}
\begin{document}
\begin{center}
{\large\bfseries 25th Irish Mathematical Olympiad\\12 May 2012}
\end{center}
\bigskip\noindent\textbf{Paper 1}\bigskip
\begin{enumerate}
\item Let
\[
C = \{1,2,3,4,5,6,7,8,9,10,11,12,13,14,15,16,17,18,19,20\}
\]
and let
\[
S = \{4,5,9,14,23,37\}.
\]
Find two sets $A$ and $B$ with the properties
\begin{enumerate}
\item $A \cap B = \emptyset$.
\item $A \cup B = C$.
\item The sum of two distinct elements of $A$ is not in $S$.
\item The sum of two distinct elements of $B$ is not in $S$.
\end{enumerate}

\item $A$, $B$, $C$ and $D$ are four points in that order on the circumference of a circle $K$. $AB$ is perpendicular to $BC$ and $BC$ is perpendicular to $CD$. $X$ is a point on the circumference of the circle between $A$ and $D$. $AX$ extended meets $CD$ extended at $E$ and $DX$ extended meets $BA$ extended at $F$.

Prove that the circumcircle of triangle $AXF$ is tangent to the circumcircle of triangle $DXE$ and that the common tangent line passes through the centre of the circle $K$.

\item Find, with proof, all polynomials $f$ such that $f$ has nonnegative integer coefficients, $f(1) = 8$ and $f(2) = 2012$.

\item There exists an infinite set of triangles with the following properties:
\begin{enumerate}
\item the lengths of the sides are integers with no common factors, and
\item one and only one angle is $60^\circ$.
\end{enumerate}
One such triangle has side lengths $5$, $7$ and $8$. Find two more.

\item
\begin{enumerate}
\item Show that if $x$ and $y$ are positive real numbers, then
\[
(x+y)^5 \;\ge\; 12xy(x^3+y^3).
\]

\item Prove that the constant $12$ is the best possible. In other words, prove that for any $K > 12$ there exist positive real numbers $x$ and $y$ such that
\[
(x+y)^5 < Kxy(x^3+y^3).
\]
\end{enumerate}
\end{enumerate}

\newpage
\begin{center}
{\large\bfseries 25th Irish Mathematical Olympiad\\12 May 2012}
\end{center}
\bigskip\noindent\textbf{Paper 2}\bigskip
\begin{enumerate}
\setcounter{enumi}{5}
\item Let $S(n)$ be the sum of the decimal digits of $n$. For example, $S(2012) = 2 + 0 + 1 + 2 = 5$. Prove that there is no integer $n > 0$ for which $n - S(n) = 9990$.

\item Consider a triangle $ABC$ with $|AB| \ne |AC|$. The angle bisector of the angle $CAB$ intersects the circumcircle of $\triangle ABC$ at two points $A$ and $D$. The circle of centre $D$ and radius $|DC|$ intersects the line $AC$ at two points $C$ and $B'$. The line $BB'$ intersects the circumcircle of $\triangle ABC$ at $B$ and $E$. Prove that $B'$ is the orthocentre of $\triangle AED$.

\item Suppose $a$, $b$, $c$ are positive numbers. Prove that
\[
\left(\frac{a}{b} + \frac{b}{c} + \frac{c}{a} + 1\right)^2 \;\ge\; (2a+b+c)\left(\frac{2}{a} + \frac{1}{b} + \frac{1}{c}\right),
\]
with equality if and only if $a = b = c$.

\item Let $x > 1$ be an integer. Prove that $x^5 + x + 1$ is divisible by at least two distinct prime numbers.

\item Let $n$ be a positive integer. A mouse sits at each corner point of an $n \times n$ square board, which is divided into unit squares as shown below for the example $n = 5$.

The mice then move according to a sequence of steps, in the following manner:
\begin{enumerate}
\item In each step, each of the four mice travels a distance of one unit in a horizontal or vertical direction. Each unit distance is called an edge of the board, and we say that each mouse uses an edge of the board.
\item An edge of the board may not be used twice in the same direction.
\item At most two mice may occupy the same point on the board at any time.
\end{enumerate}
The mice wish to collectively organise their movements so that each edge of the board will be used twice (not necessarily by the same mouse), and each mouse will finish up at its starting point. Determine, with proof, the values of $n$ for which the mice may achieve this goal.
\end{enumerate}
\end{document}
